% !TeX program = pdflatex
\documentclass[11pt,a4paper]{article}
\usepackage[T1]{fontenc}
\usepackage[utf8]{inputenc}
\usepackage[polish]{babel}
\usepackage{lmodern,amsmath,amssymb,amsthm,geometry,microtype,hyperref}
\geometry{margin=2.5cm}
\setlength{\parindent}{0pt}
\setlength{\parskip}{5pt}
\newcommand{\Hom}{\operatorname{Hom}}
\newcommand{\id}{\mathrm{id}}
\newcommand{\Set}{\mathsf{Set}}
\newcommand{\Rel}{\mathsf{Rel}}
\newcommand{\Vect}{\mathsf{Vect}}
\newcommand{\Grp}{\mathsf{Grp}}
\title{Teoria kategorii — zbiór zadań\\\large Zestaw podstawowy z podpowiedziami}
\author{}
\date{}

% BEGIN GLOBAL CODE LISTINGS (generated from latex/pl/code-listings.tex)
\usepackage{xcolor}
\usepackage{listings}
\definecolor{codebg}{RGB}{248,249,251}
\definecolor{codeborder}{RGB}{180,191,203}
\definecolor{codecomment}{RGB}{61,113,79}
\definecolor{codekeyword}{RGB}{41,78,132}
\lstdefinelanguage{Julia}{
  morekeywords={abstract,baremodule,begin,break,catch,const,continue,do,else,elseif,
    end,export,false,finally,for,function,global,if,import,in,isa,let,local,macro,
    module,mutable,primitive,quote,return,struct,true,try,type,using,where,while,
    missing,nothing,println,print},
  sensitive=true,morecomment=[l]{\#},morecomment=[n]{\#=}{=\#},
  morestring=[b]",
}
\lstdefinestyle{skrypt}{
  backgroundcolor=\color{codebg},frame=single,rulecolor=\color{codeborder},
  framesep=5pt,frameround=tttt,framerule=.5pt,
  basicstyle=\ttfamily\footnotesize,keywordstyle=\color{codekeyword}\bfseries,
  commentstyle=\color{codecomment}\itshape,stringstyle=\color{purple!65!black},
  numbers=left,numberstyle=\tiny\color{gray!75!black},numbersep=8pt,
  stepnumber=1,breaklines=true,breakatwhitespace=false,
  columns=fullflexible,keepspaces=true,showstringspaces=false,tabsize=4,
  linewidth=.90\linewidth,captionpos=b,aboveskip=11pt,belowskip=12pt,
  xleftmargin=12pt,xrightmargin=2pt,
}
\lstdefinestyle{pythonpde}{style=skrypt,language=Python}
\lstdefinestyle{python}{style=skrypt,language=Python}
\lstdefinestyle{julia}{style=skrypt,language=Julia}
\lstset{style=skrypt}
\renewcommand{\lstlistingname}{Kod}
\renewcommand{\lstlistlistingname}{Spis kodów}
% END GLOBAL CODE LISTINGS
\begin{document}
\maketitle
Zestaw towarzyszy rozdziałom CAT-00--18 skryptu \emph{Teoria kategorii od podstaw}.
Identyfikatory CAT-Z-00--18 wskazują odpowiednie rozdziały źródłowe, a nie numery sekcji PDF.
W każdym zadaniu należy uzasadnić istnienie, dobre określenie i jednoznaczność konstruowanych morfizmów, gdy są wymagane. Podpowiedzi znajdują się w drugiej części; nie są pełnymi rozwiązaniami.
\section{Zadania}
\begin{enumerate}
\item[\textbf{CAT-Z-00}] W kategorii zbiorów i funkcji injektywnych wykaż, że nie istnieje obiekt końcowy. Ustal, czy istnieje obiekt początkowy.
\item[\textbf{CAT-Z-01}] Dla zbioru $S$ zdefiniuj kategorię, w której morfizmem $x\to y$ jest dokładnie jedna para $(x,y)$ dla każdego $x,y\in S$. Sprawdź aksjomaty kategorii i wykaż, że jest grupoidem. Uwzględnij $S=\varnothing$.
\item[\textbf{CAT-Z-02}] W kategorii zbiorów częściowo uporządkowanych i funkcji monotonicznych podaj bijekcję, która nie jest izomorfizmem. Wykaż bezpośrednio, że jest jednocześnie monomorfizmem i epimorfizmem.
\item[\textbf{CAT-Z-03}] Wykaż, że przyporządkowanie $\mathcal P(X)$ wraz z przeciwobrazem jest funktorem kontrawariantnym, a wraz z obrazem bezpośrednim kowariantnym. Czy oba funktory są wierne?
\item[\textbf{CAT-Z-04}] Ustal wszystkie transformacje naturalne funktora identycznościowego na $\Grp$ w siebie. Wykorzystaj grupę wolną na jednym i na dwóch generatorach.
\item[\textbf{CAT-Z-05}] Wyznacz obiekt początkowy i końcowy w kategorii przestrzeni topologicznych punktowanych. Uzasadnij, że morfizm zerowy jest odwzorowaniem stałym o wartości w punkcie wyróżnionym przeciwdziedziny.
\item[\textbf{CAT-Z-06}] Wykaż własność uniwersalną iloczynu dowolnej rodziny przestrzeni wektorowych nad $k$ oraz koproduktu tej rodziny. Dla nieskończonej rodziny kopii $k$ wskaż element iloczynu, który nie należy do sumy prostej.
\item[\textbf{CAT-Z-07}] Niech $f:A\to C$ i $g:B\to C$ będą odwzorowaniami liniowymi. Zbuduj pullback w $\Vect_k$ i udowodnij jego własność uniwersalną. Kiedy rzut na $B$ jest surjektywny?
\item[\textbf{CAT-Z-08}] Dla układu $X_n=\mathbb Z/2^n\mathbb Z$, $n\geq1$, z naturalnymi rzutami $X_{n+1}\to X_n$ opisz granicę odwrotną w $\Set$. Uzasadnij bijekcję między nią a zbiorem nieskończonych ciągów cyfr binarnych.
\item[\textbf{CAT-Z-09}] Dla ustalonego ciała $k$ i zbioru $S$ wykaż reprezentowalność funktora $\Vect_k\to\Set$, $V\mapsto\Set(S,V)$, gdzie $V$ po prawej oznacza zbiór wektorów.
\item[\textbf{CAT-Z-10}] Sklasyfikuj wszystkie transformacje naturalne $X^n\Rightarrow X^m$ w $\Set$ dla ustalonych liczb naturalnych $n,m$, włącznie z zerem.
\item[\textbf{CAT-Z-11}] Wyprowadź z własności uniwersalnej obiektu wykładniczego kanoniczny izomorfizm $(C^B)^A\cong C^{A\times B}$ w kategorii kartezjańsko domkniętej.
\item[\textbf{CAT-Z-12}] Ustal, czym jest iloczyn dwóch obiektów w kategorii $\Set/B$. Zbadaj przypadek, gdy $B=\varnothing$.
\item[\textbf{CAT-Z-13}] Dla ustalonego ciała $k$ zdefiniuj kategorię o obiektach $n\geq0$ i morfizmach $n\to m$ danych przez macierze $m\times n$. Wykaż jej równoważność z kategorią skończenie wymiarowych przestrzeni wektorowych (w ustalonym uniwersum).
\item[\textbf{CAT-Z-14}] Dla funkcji $f:X\to Y$ i zbiorów potęgowych uporządkowanych przez inkluzję udowodnij oba sprzężenia $\exists_f\dashv f^{-1}\dashv\forall_f$. Zbadaj zachowanie $\forall_f$ na punktach spoza $f(X)$.
\item[\textbf{CAT-Z-15}] Dla monady list na $\Set$ udowodnij, że morfizmy jej algebr są dokładnie homomorfizmami odpowiednich monoidów. Uwzględnij listę pustą.
\item[\textbf{CAT-Z-16}] W kategorii przestrzeni unormowanych i ograniczonych odwzorowań liniowych udowodnij, że każdy morfizm o gęstym obrazie jest epimorfizmem. Wykaż odwrotność przy pomocy ilorazu przez domknięcie obrazu.
\item[\textbf{CAT-Z-17}] Udowodnij, że suma rozłączna $A\sqcup B$ jest jednocześnie iloczynem i koproduktem w $\Rel$. Wykaż, że dla $A=B=\{*\}$ zwykły iloczyn kartezjański nie może pełnić roli iloczynu w $\Rel$.
\item[\textbf{CAT-Z-18}] Rozważ sprzężenie funktorów $F:\mathcal C\to\mathcal D$ i $G:\mathcal D\to\mathcal C$. Podaj typy jednostki, kojednostki i mnożenia monady $T=GF$. Zapisz obie tożsamości trójkątne i prawa jednostki monady, wskazując, które tożsamości trójkątne je uzasadniają.
\end{enumerate}
\newpage
\section{Podpowiedzi}
\begin{enumerate}
\item[\textbf{CAT-Z-00}] Dla kandydata $T$ na obiekt końcowy rozważ zbiór jednoelementowy. Jedyność injekcji wymusza, że $T$ jest jednoelementowy; następnie rozważ zbiór dwuelementowy. Z pustego zbioru istnieje dokładnie jedna funkcja do każdego zbioru.
\item[\textbf{CAT-Z-01}] Złożenie $(y,z)\circ(x,y)=(x,z)$; identyczność $(x,x)$; odwrotność $(y,x)$. Dla pustego $S$ aksjomaty są spełnione w sposób pusty.
\item[\textbf{CAT-Z-02}] Użyj identyczności zbioru dwuelementowego z porządkiem dyskretnym do łańcucha. Odwrotność nie zachowuje porządku. Skracanie wynika z własności funkcji.
\item[\textbf{CAT-Z-03}] Zapisz wzory na obrazy i przeciwobrazy złożeń. Obraz singletonu rozpoznaje wartość funkcji, a przeciwobrazy singletonów rozpoznają jej włókna.
\item[\textbf{CAT-Z-04}] Na grupie wolnej $\mathbb Z$ składowa ma postać $n\mapsto rn$. Naturalność względem homomorfizmu $\mathbb Z\to G$ wymusza $\tau_G(g)=g^r$. W grupie wolnej na $x,y$ warunek $(xy)^r=x^ry^r$ pozostawia tylko $r=0,1$: porównaj zredukowane słowa, osobno dla $r<0$ i $r>1$.
\item[\textbf{CAT-Z-05}] Obiekt jednoelementowy z jedynym punktem. Każde odwzorowanie z niego zachowujące punkt jest wymuszone, podobnie odwzorowanie do niego.
\item[\textbf{CAT-Z-06}] Do iloczynu kieruj morfizm współrzędnie. Z sumy prostej wysyłaj wektor na skończoną sumę obrazów jego współrzędnych. Rozważ rodzinę samych jedynek.
\item[\textbf{CAT-Z-07}] $P=\{(a,b):f(a)=g(b)\}=\ker(f,-g)$. Rzut na $B$ jest surjektywny dokładnie wtedy, gdy $g(B)\subseteq f(A)$.
\item[\textbf{CAT-Z-08}] W granicy znajdują się zgodne rodziny klas reszt. Każda klasa modulo $2^n$ ma reprezentanta z przedziału $[0,2^n-1]$, a zgodność wymusza zgodność pierwszych $n$ cyfr zapisu binarnego liczonego od najmniej znaczącej.
\item[\textbf{CAT-Z-09}] Obiekt reprezentujący to $k^{(S)}$. Funkcja na generatorach rozszerza się jednoznacznie do odwzorowania liniowego. Naturalność sprawdź przez składanie.
\item[\textbf{CAT-Z-10}] $X^n\cong\Hom_{\Set}(\{1,\ldots,n\},X)$. Zastosuj kowariantną postać lematu Yonedy; naturalna operacja odpowiada funkcji ze zbioru $m$-elementowego do $n$-elementowego. W szczególności dla $n=0<m$ nie ma żadnej, a dla $m=0$ jest jedna.
\item[\textbf{CAT-Z-11}] Dla dowolnego $X$ porównaj naturalnie $\Hom(X,(C^B)^A)$ z $\Hom((X\times A)\times B,C)$, następnie użyj kanonicznej łączności iloczynu i lematu Yonedy.
\item[\textbf{CAT-Z-12}] Zbiór par o równych obrazach w $B$ z ich wspólnym odwzorowaniem do $B$. Dla pustego $B$ jedyny obiekt to odwzorowanie z pustego zbioru do niego.
\item[\textbf{CAT-Z-13}] Funktor wysyła $n$ na $k^n$, a macierz na odwzorowanie liniowe zadane przez mnożenie macierzy przez wektor. Jest w pełni wierny i istotnie surjektywny. Wskaż konwencję wyboru baz.
\item[\textbf{CAT-Z-14}] $\exists_f(A)=f(A)$ oraz $\forall_f(A)=\{y:f^{-1}(\{y\})\subseteq A\}$. Porównaj inkluzje z przeciwobrazem. Puste włókno jest podzbiorem każdego $A$.
\item[\textbf{CAT-Z-15}] Składowa algebry $a$ daje jednostkę $a([])$ oraz mnożenie $a([x,y])$. Dla morfizmu algebr $h$ zastosuj $ha=bT(h)$ do list pustej i dwuelementowej. Odwrotnie użyj indukcji po długości listy.
\item[\textbf{CAT-Z-16}] Dwa ciągłe odwzorowania równe na gęstym podzbiorze są równe wszędzie, ponieważ przeciwdziedzina jest Hausdorffa. Jeśli domknięcie obrazu jest właściwe, rzut na niezerową przestrzeń ilorazową i morfizm zerowy są różne, ale mają równe złożenia z badanym morfizmem.
\item[\textbf{CAT-Z-17}] Użyj grafów włączeń i ich relacji odwrotnych. Relacja do sumy rozłącznej jest jednoznacznie wyznaczona przez ograniczenia przeciwdziedziny do składników. Dla jednopunktowego obiektu testowego porównaj liczby relacji: do singletonu są dwie, a par relacji do dwóch singletonów cztery.
\item[\textbf{CAT-Z-18}] $\eta:\id_{\mathcal C}\Rightarrow GF$, $\varepsilon:FG\Rightarrow\id_{\mathcal D}$, $\mu_A=G(\varepsilon_{F(A)})$. Rozpisz $\mu_A T(\eta_A)$ i $\mu_A\eta_{T(A)}$; skorzystaj odpowiednio z trójkątów dla $F$ i $G$.
\end{enumerate}
\end{document}
